%% Question bank chapter file - generated by bank_tools.py build (do not edit while the bank is still being built:
%% it is regenerated from the per-test data in Question Bank\data). Every question is one bankq block:
%%   \begin{bankq}[id=..., source={...}, paper=.., qno=.., type=..., subtopic={...}, answer={...}, answer_src=key|solved,
%%                 working={...}, difficulty=1-10, diff_parts={K S I C F}, diff_note={...},
%%                 group=<passage id or empty>, also_in={duplicate copies in other tests}]
%% difficulty = min(10, K+S+I+C+F-1): K concepts 1-3, S steps 1-4, I insight 0-2, C computation 0-2, F format 0-1
%%   ...question LaTeX, options, figures (\bfig / \bfigi -> img/)...
%%   \end{bankq}
%% Passages shared by several questions are bankpassage blocks; their questions carry group=<passage id>.
\documentclass[10pt]{article}
\usepackage[a4paper,margin=1.8cm]{geometry}
\usepackage[T1]{fontenc}
\usepackage{amsmath,amssymb,mathtools}
\usepackage{graphicx}
\usepackage[export]{adjustbox}
\usepackage{xcolor,needspace,keyval}
\usepackage{enumitem}
\usepackage{booktabs,array,tabularx}
\usepackage[hidelinks]{hyperref}
\setlength{\parindent}{0pt}
\setlength{\parskip}{3pt}
\makeatletter
\define@key{bq}{id}{\def\bq@id{#1}}\define@key{bq}{source}{\def\bq@src{#1}}\define@key{bq}{paper}{\def\bq@paper{#1}}
\define@key{bq}{qno}{\def\bq@qno{#1}}\define@key{bq}{type}{\def\bq@type{#1}}\define@key{bq}{subtopic}{\def\bq@sub{#1}}
\define@key{bq}{answer}{\def\bq@ans{#1}}\define@key{bq}{answer_src}{\def\bq@asrc{#1}}\define@key{bq}{working}{\def\bq@work{#1}}
\define@key{bq}{group}{\def\bq@group{#1}}\define@key{bq}{also_in}{\def\bq@also{#1}}
\define@key{bq}{difficulty}{\def\bq@diff{#1}}\define@key{bq}{diff_parts}{\def\bq@dparts{#1}}\define@key{bq}{diff_note}{\def\bq@dnote{#1}}
\newenvironment{bankq}[1][]{\def\bq@work{}\def\bq@also{}\def\bq@group{}\def\bq@diff{?}\def\bq@dparts{}\def\bq@dnote{}\setkeys{bq}{#1}%
  \par\needspace{4\baselineskip}\medskip\noindent{\small\textbf{\bq@id}\hfill\color{gray}[\bq@src, Paper \bq@paper, Q\bq@qno]
  \bq@type\ $\cdot$ \bq@sub\ $\cdot$ \textbf{D\bq@diff}}\par\nopagebreak\relax}
  {\par{\small\color{teal}Answer: \textbf{\bq@ans}\ifx\bq@work\@empty\else\ (solved: \bq@work)\fi
  \ifx\bq@also\@empty\else\par\color{gray}Also in: \bq@also\fi}\par}
\define@key{bp}{id}{\def\bp@id{#1}}\define@key{bp}{source}{\def\bp@src{#1}}\define@key{bp}{questions}{\def\bp@qs{#1}}
\newenvironment{bankpassage}[1][]{\setkeys{bp}{#1}\par\needspace{4\baselineskip}\bigskip\noindent
  {\small\textbf{Passage \bp@id}\hfill\color{gray}[\bp@src; questions \bp@qs]}\par\nopagebreak\relax\itshape}{\par}
\makeatother
\newcommand{\bfig}[2]{\par\smallskip{\centering\includegraphics[width=#1pt,max width=0.62\linewidth,max height=0.32\textheight]{img/#2}\par}\smallskip}
\newcommand{\bfigi}[2]{\includegraphics[width=#1pt,max width=\linewidth,max height=0.32\textheight,valign=c]{img/#2}}
\newcommand{\bmissing}[1]{\textit{[figure missing in the source conversion: \detokenize{#1}]}}
\newenvironment{qoptions}{\begin{itemize}[leftmargin=1.8em,itemsep=1pt,topsep=2pt,label={}]}{\end{itemize}}
\newcommand{\opt}[2]{\item[(#1)] #2}
\providecommand{\fwstart}{}\providecommand{\fwend}{}\providecommand{\qfullsplit}{}
\newcolumntype{L}[1]{>{\raggedright\arraybackslash}p{#1}}
\newcolumntype{C}[1]{>{\centering\arraybackslash}p{#1}}
\newcolumntype{R}[1]{>{\raggedleft\arraybackslash}p{#1}}
\providecommand{\sectionhdr}[1]{}\providecommand{\subsectionhdr}[1]{}


\begin{document}
{\Large\bfseries Probability I + II --- Question Bank}\par\smallskip
{\small Mathematics chapter M06. Questions from all converted test papers, grouped by source series. Answers: printed key where it matches the paper, otherwise solved (answer\_src=solved, with working). Exact repeats across papers appear once, with also\_in listing the other copies.}\par\medskip\hrule

\section*{Score-I papers}
\begin{bankq}[id=SI12-P1-MATH-S2-Q04, source={Score-I - Test 12 (09 Mar 2022)}, paper=1, qno=4, type=numerical, subtopic={Classical probability}, answer={20.00}, answer_src=key, difficulty=5, diff_parts={K2 S2 I1 C0 F1}, diff_note={concepts: arrangements with parity constraints, stars and bars | steps: girls at the ends; odd boys in each of 4 gaps b\_i = 2c\_i + 1; sum c\_i = 3 gives 20; total C(15,5) = 3003 | I: transform the odd condition}]
5 girls and 10 boys sit at random in a row having 15 chairs numbered as 1 to 15. If the probability that the end seats are occupied by the girls and between any two girls odd number of boys take seat is $\dfrac n{3003}$, then find the value of $n$
\end{bankq}
\begin{bankq}[id=SI12-P1-MATH-S2-Q05, source={Score-I - Test 12 (09 Mar 2022)}, paper=1, qno=5, type=numerical, subtopic={Bayes}, answer={48.78}, answer_src=key, difficulty=3, diff_parts={K1 S2 I0 C0 F1}, diff_note={concepts: Bayes' theorem | steps: probability of white for each discard case; total probability; Bayes}]
A bag contains 5 white and 3 red balls. A ball is selected at random and discarded after noting its colour and, two balls of the other colour put in the bag. A ball now drawn from the bag is white and the probability that discarded ball was also white is $K$, then find $100K$.
\end{bankq}
\begin{bankpassage}[id=SI12-P2-MATH-S1-G3-4, source={Score-I - Test 12 (09 Mar 2022)}, questions={3-4}]
Three alike dice are rolled simultaneously.\\
Column 1 : Sum of numbers appeared on the dice are given.\\
Column 2 : Number of ways in which the event can occur.\\
Column 3 : Probability of the event.
\begin{center}\footnotesize
\begin{tabular}{|p{0.6cm}|p{2.2cm}|p{0.6cm}|p{0.8cm}|p{0.5cm}|p{1.3cm}|}
\hline
 & \centering Column 1 & \multicolumn{2}{c|}{Column 2} & \multicolumn{2}{c|}{Column 3}\\
\hline
(I) & 6 & (i) & 3 & (P) & $\dfrac12$\\
(II) & 15 & (ii) & 10 & (Q) & $\dfrac5{108}$\\
(III) & less than 11 & (iii) & 28 & (R) & $\dfrac3{56}$\\
(IV) & greater than or equal to 11 & (iv) & 108 & (S) & $\dfrac5{28}$\\
\hline
\end{tabular}
\end{center}
\end{bankpassage}
\begin{bankq}[id=SI12-P2-MATH-S1-Q03, source={Score-I - Test 12 (09 Mar 2022)}, paper=2, qno=3, type=single, subtopic={Classical probability}, answer={B}, answer_src=key, difficulty=3, diff_parts={K1 S2 I0 C1 F0}, diff_note={concepts: classical probability | steps: count for item I; count for item II; count for item III; count for item IV and match}, group=SI12-P2-MATH-S1-G3-4]
Identify the correct combination ?
\begin{qoptions}
\opt{A}{(I) (ii) (Q)}
\opt{B}{(I) (i) (Q)}
\opt{C}{(II) (i) (R)}
\opt{D}{(IV) (iv) (P)}
\end{qoptions}
\end{bankq}
\begin{bankq}[id=SI12-P2-MATH-S1-Q04, source={Score-I - Test 12 (09 Mar 2022)}, paper=2, qno=4, type=single, subtopic={Classical probability}, answer={C}, answer_src=key, difficulty=3, diff_parts={K1 S2 I0 C1 F0}, diff_note={concepts: classical probability | steps: count for item I; count for item II; count for item III; count for item IV and match}, group=SI12-P2-MATH-S1-G3-4]
Identify the correct combination ?
\begin{qoptions}
\opt{A}{(I) (i) (R)}
\opt{B}{(II) (ii) (R)}
\opt{C}{(III) (iii) (P)}
\opt{D}{(IV) (iv) (P)}
\end{qoptions}
\end{bankq}

\section*{Score Advanced papers}
\begin{bankq}[id=SA01-P1-MATH-S3-Q03, source={Score Advanced - Test 01 (22 Sep 2020)}, paper=1, qno=3, type=numerical, subtopic={Conditional probability}, answer={0.315}, answer_src=key, difficulty=8, diff_parts={K2 S2 I2 C2 F1}, diff_note={concepts: probability with cards, conditional probability | steps: case of one honour card; conditional probability; value}]
A pack of playing card is consist of 51 cards. Cards are drawn one by one from the pack without replacement. If first 4 cards drawn consist of only one honour card, then the probability that the missing card being honour card is $\dfrac ab$ (where $a$ and $b$ are co-prime and all Jacks, Queens, Kings and Aces are considered as honour card) then $\dfrac ab$ is\underline{\hspace{1.5cm}}.
\end{bankq}
\begin{bankq}[id=SA01-P1-MATH-S3-Q04, source={Score Advanced - Test 01 (22 Sep 2020)}, paper=1, qno=4, type=numerical, subtopic={Classical probability}, answer={2.50}, answer_src=key, difficulty=8, diff_parts={K2 S2 I2 C2 F1}, diff_note={concepts: extinction probability of a branching process, classical probability | steps: equation for the probability; solve; choose the root}]
A certain kind of bacteria either die, split into two or split into three bacteria. All splits are exact copies. The chances of dying is $\dfrac14$, the chances of splitting into two is $\dfrac12$ and splitting into three is $\dfrac14$. If the probability that it survives for infinite length of time is $\dfrac{\mathrm m-\sqrt{13}}{\mathrm n}$ (m,n $\in$ N), then the value of (m/n) is
\end{bankq}
\begin{bankq}[id=SA01-P2-MATH-S3-Q08, source={Score Advanced - Test 01 (22 Sep 2020)}, paper=2, qno=8, type=numerical, subtopic={Conditional probability}, answer={29.00}, answer_src=key, difficulty=7, diff_parts={K2 S2 I2 C1 F1}, diff_note={concepts: probability with a revolver, conditional probability | steps: cases; conditional probability; value}]
A game of luck is designed in which you have to shoot a target. The organizer puts two bullet in an empty six round revolver and spins. You point it towards balloon and shoot. It was an empty shot (i.e. bullet was not fired). If the bullets are placed randomly then the probability that you will shoot the balloon in next shot is $\dfrac PQ$ (Where P and Q are relatively coprimes) \{It is assumed that fired bullet always hits the target\}, then the value of (P$^2$ + Q$^2$) is
\end{bankq}
\begin{bankq}[id=SA02-P1-MATH-S1-Q05, source={Score Advanced - Test 02 (13 May 2021)}, paper=1, qno=5, type=multi, subtopic={Classical probability}, answer={A,B}, answer_src=solved, working={Midpoint uniform in the disk: probability 5/16}, difficulty=7, diff_parts={K2 S2 I2 C1 F1}, diff_note={concepts: probability for a random chord, classical probability | steps: model for the chord; relation of length and distance; probability}]
Probability that the length of a randomly chosen chord of a circle lies between $\dfrac{1}{2}$ and $\dfrac{3}{4}$ of its diameter is $\dfrac{p}{q}$ where p and q are coprime number, then
\begin{qoptions}
\opt{A}{p + q = 21}
\opt{B}{q $-$ p = 11}
\opt{C}{p + q = 20}
\opt{D}{q $-$ p = 10}
\end{qoptions}
\end{bankq}
\begin{bankq}[id=SA02-P2-MATH-S1-Q03, source={Score Advanced - Test 02 (13 May 2021)}, paper=2, qno=3, type=multi, subtopic={Binomial distribution}, answer={A,C,D}, answer_src=solved, working={Probability C(2n,n)/4\^{}n; the product with (n+r)/(2r) gives C(2n,n)/2\^{}n}, difficulty=7, diff_parts={K2 S2 I2 C1 F1}, diff_note={concepts: probability of equal heads, binomial distribution | steps: binomial coefficient; products; check}]
Probability of n heads in 2n tosses of a fair coin can be given by
\begin{qoptions}
\opt{A}{$\displaystyle\prod_{r=1}^{n}\left(\dfrac{2r - 1}{2r}\right)$}
\opt{B}{$\displaystyle\prod_{r=1}^{n}\left(\dfrac{n + r}{2r}\right)$}
\opt{C}{$\displaystyle\sum_{r=0}^{n}\left(\dfrac{{}^nC_r}{2^n}\right)^2$}
\opt{D}{$\dfrac{\displaystyle\sum_{r=0}^{n}\left({}^nC_r\right)^2}{\left(\displaystyle\sum_{r=0}^{2n}{}^{2n}C_r\right)}$}
\end{qoptions}
\end{bankq}
\begin{bankq}[id=SA05-P1-MATH-S1-Q14, source={Score Advanced - Test 05 (15 Jul 2021)}, paper=1, qno=14, type=single, subtopic={Binomial distribution}, answer={D}, answer_src=key, difficulty=4, diff_parts={K2 S1 I1 C1 F0}, diff_note={concepts: binomial distribution from mean and variance and np=4 and npq=4/3 and so q=1/3, n=6 | steps: np=4, npq=4/3, so q=1/3, n=6; P(X>=2)=1-P(0)-P(1)}]
Mean and variance of binomial probability distribution are 4 and $\dfrac{4}{3}$ respectively then probability of atleast two success is equal to
\begin{qoptions}
\opt{A}{$\dfrac{552}{729}$}
\opt{B}{$\dfrac{201}{243}$}
\opt{C}{$\dfrac{298}{343}$}
\opt{D}{$\dfrac{716}{729}$}
\end{qoptions}
\end{bankq}
\begin{bankq}[id=SA06-P1-MATH-S3-Q05, source={Score Advanced - Test 06 (04 Sep 2021)}, paper=1, qno=5, type=numerical, subtopic={Classical probability}, answer={6}, answer_src=key, difficulty=8, diff_parts={K3 S1 I2 C2 F1}, diff_note={concepts: probability with random removal of cards, classical probability, condition on the number of aces removed | steps: condition on the number of aces removed}]
All the face cards from a pack of 52 playing cards are removed. From the remaining pack half of the cards are randomly removed whithout looking at them and then randomly drawn two cards simultaneously from the remaining. If the probability that, two cards drawn are both aces, is $\dfrac{p\left({}^{38}C_{20}\right)}{{}^{40}C_{20}\cdot{}^{20}C_2}$, then the value of p is ?
\end{bankq}
\begin{bankq}[id=SA06-P2-MATH-S3-Q05, source={Score Advanced - Test 06 (04 Sep 2021)}, paper=2, qno=5, type=numerical, subtopic={Classical probability}, answer={4}, answer_src=key, difficulty=8, diff_parts={K3 S1 I2 C2 F1}, diff_note={concepts: probability with a hypergeometric structure, classical probability, Printed key | steps: Printed key}]
A circle of radius r = 4 units is inscribed in an equilateral triantgle ABC, then an equilateral triangle is inscribed in the circle, a circle again is inscribed in the later triangle and so on. In this way the process continues infinitely. If r, x$_1$, x$_2$, ..., x$_n$, .... be the radii of these circles respectively, then the sum of radii of all the circles x$_1$ + x$_2$ + x$_3$ ...$\infty$ is equals to
\end{bankq}
\begin{bankq}[id=SA07-P1-MATH-S1-Q02, source={Score Advanced - Test 07 (16 Apr 2023)}, paper=1, qno=2, type=multi, subtopic={Conditional probability}, answer={A,B}, answer_src=key, difficulty=8, diff_parts={K3 S1 I2 C2 F1}, diff_note={concepts: dice sequence not greater than the preceding number, conditional probability, count non-increasing sequences with a single 1 in the fourth place | steps: count non-increasing sequences with a single 1 in the fourth place}]
A fair dice is rolled four times. Find the probability that each number is not greater than the preceding number. When it is known that it shows `1' only once \& that too in the fourth round.
\begin{qoptions}
\opt{A}{$\dfrac{7}{25}$}
\opt{B}{less than $\dfrac{1}{3}$}
\opt{C}{greater than $\dfrac{1}{3}$}
\opt{D}{$\dfrac{35}{216}$}
\end{qoptions}
\end{bankq}
\begin{bankq}[id=SA07-P2-MATH-S2-Q09, source={Score Advanced - Test 07 (16 Apr 2023)}, paper=2, qno=9, type=single, subtopic={Bayes}, answer={D}, answer_src=key, difficulty=8, diff_parts={K3 S2 I2 C2 F0}, diff_note={concepts: posterior probabilities with three coins, bayes, likelihoods of two heads in three tosses | steps: likelihoods of two heads in three tosses; Bayes; fourth toss}]
A man has 3 coins A, B and C, A is unbiased, the probability that a head will result when B is tossed is 2/3, the probability that a head will result when C is tossed is 1/3. One of the coin chosen at random and is tossed 3 times, giving a total of two heads and one tail.\par
{\scriptsize\setlength{\tabcolsep}{1.5pt}\noindent\begin{tabular}{|L{.6cm}|L{5.2cm}|L{.6cm}|L{1.2cm}|}
\hline
 & \multicolumn{1}{c|}{List-I} & & List-II \\ \hline
(P) & The probability that the chosen coin was A & (1) & 32/75 \\ \hline
(Q) & The probability that the chosen coin was B & (2) & 1/2 \\ \hline
(R) & The probability that the fourth toss will result in head & (3) & 209/1296 \\ \hline
(S) & The probability that the fourth toss will result in tail & (4) & 9/25 \\ \hline
 & & (5) & 215/1296 \\ \hline
\end{tabular}}\par
\begin{qoptions}
\opt{A}{P $\rightarrow$ 4;Q $\rightarrow$ 1;R $\rightarrow$ 2;S $\rightarrow$ 3}
\opt{B}{P $\rightarrow$ 5;Q $\rightarrow$ 4;R $\rightarrow$ 2;S $\rightarrow$ 3}
\opt{C}{P $\rightarrow$ 3;Q $\rightarrow$ 2;R $\rightarrow$ 5;S $\rightarrow$ 1}
\opt{D}{P $\rightarrow$ 4;Q $\rightarrow$ 1;R $\rightarrow$ 2;S $\rightarrow$ 2}
\end{qoptions}
\end{bankq}
\begin{bankq}[id=SA07-P2-MATH-S3-Q06, source={Score Advanced - Test 07 (16 Apr 2023)}, paper=2, qno=6, type=single, subtopic={Conditional probability}, answer={5}, answer_src=key, difficulty=7, diff_parts={K3 S1 I2 C2 F0}, diff_note={concepts: posterior probability that die A was used, conditional probability, Bayes for n red throws | steps: Bayes for n red throws}]
Die A has 2 white and 4 red faces whereas die B has 4 white and 2 red faces. A coin is flipped once. If it shows a head, the game continues by throwing die A, if it shows tail, then die B is to be used. If the probability that die A is used is $\dfrac{32}{33}$ when it is given that red turns up every time in first n throws. Then find the value of n \underline{\hspace{1.2cm}}.
\end{bankq}

\section*{TOAS papers}
\begin{bankq}[id=TOAS18-P1-MATH-S1-Q04, source={TOAS - Test 18 (30 Oct 2022)}, paper=1, qno=4, type=multi, subtopic={Classical probability}, answer={C,D}, answer_src=solved, working={$P(E_1)=9/25$, $P(E_2)=9/100$; $E_2\subset E_1$ so $P(E_2/E_1)=1/4$ and $P(E_1/E_2)=1$}, difficulty=3, diff_parts={K1 S2 I0 C0 F1}, diff_note={concepts: probability of unit digits in products | steps: probability of E1; probability of E2; conditional probabilities}]
Two whole numbers are randomly selected and multiplied. Consider two events $E_1$ and $E_2$ defined as
\par\noindent$E_1$ : Their product is divisible by 5\par\noindent$E_2$ : Unit's place in their product is 5.\par
Which of the following statement(s) is/are correct ?
\begin{qoptions}
\opt{A}{$E_1$ is twice as likely to occur as $E_2$}
\opt{B}{$E_1$ and $E_2$ are disjoint}
\opt{C}{$P(E_2/E_1)=1/4$}
\opt{D}{$P(E_1/E_2)=1$}
\end{qoptions}
\end{bankq}
\begin{bankq}[id=TOAS18-P1-MATH-S2-Q09, source={TOAS - Test 18 (30 Oct 2022)}, paper=1, qno=9, type=single, subtopic={Classical probability}, answer={D}, answer_src=solved, working={Total $\binom{20}{4}=4845$: no pair $3360/4845=\tfrac{14\times16}{19\times17}$; at least one pair $=\tfrac{\binom{20}4-\binom{10}42^4}{\binom{20}4}$; exactly one pair $\tfrac{96}{19\times17}$; two pairs $\tfrac{3}{19\times17}$}, difficulty=2, diff_parts={K1 S2 I0 C0 F0}, diff_note={concepts: probability of matching shoes | steps: total ways; counts for no pair, one pair, two pairs; match the fractions}]
A showroom has 10 different pairs of expensive shoes. A thief steals 4 shoes randomly from the showroom. Match the following :
\begin{center}\footnotesize\setlength{\tabcolsep}{1pt}
\begin{tabular}{|L{0.6cm}L{3.4cm}|L{0.6cm}L{2.4cm}|}\hline
&List-I&&List-II\\\hline
(I)&The probability that there is no matching pair among the 4 shoes stolen is&(P)&$\dfrac3{19\times17}$\\\hline
(II)&The probability that there is atleast one matching pair among the 4 shoes stolen is&(Q)&$\dfrac{96}{19\times17}$\\\hline
(III)&The probability that there is exactly one matching pair among the 4 shoes stolen is&(R)&$\dfrac{{}^{20}C_4-{}^{10}C_4\cdot2^4}{{}^{20}C_4}$\\\hline
(IV)&The probability that there are two matching pair among the 4 shoes stolen is&(S)&$\dfrac{14\times16}{19\times17}$\\\hline
&&(T)&$\dfrac{16}{19\times17}$\\\hline
\end{tabular}\end{center}
\begin{qoptions}
\opt{A}{I $\to$ R;II $\to$ S;III $\to$ P;IV $\to$ T}
\opt{B}{I $\to$ S;II $\to$ R;III $\to$ T;IV $\to$ Q}
\opt{C}{I $\to$ R;II $\to$ S;III $\to$ Q;IV $\to$ P}
\opt{D}{I $\to$ S;II $\to$ R;III $\to$ Q;IV $\to$ P}
\end{qoptions}
\end{bankq}
\begin{bankq}[id=TOAS18-P1-MATH-S3-Q07, source={TOAS - Test 18 (30 Oct 2022)}, paper=1, qno=7, type=numerical, subtopic={Classical probability}, answer={6}, answer_src=solved, working={15 equally likely non-empty answer sets, success chance $1/15$ per attempt: $1-(14/15)^n\ge1/3$ gives $n\ge5.88$, so 6}, difficulty=4, diff_parts={K1 S2 I1 C0 F1}, diff_note={concepts: probability of a correct subset | steps: number of non-empty subsets; probability per attempt; inequality for the number of chances}]
In a multiple choice question there are four alternative answers of which one or more than one is correct. A candidate will get marks on the question only if he ticks all the correct answer. The candidate ticks the answers at random. The least number of chances he should be given if the probability of getting marks is greater than or equal to $\dfrac13$ is -
\end{bankq}
\begin{bankq}[id=TOAS18-P2-MATH-S2-Q08, source={TOAS - Test 18 (30 Oct 2022)}, paper=2, qno=8, type=multi, subtopic={Conditional probability}, answer={A,B,C,D}, answer_src=key, difficulty=2, diff_parts={K1 S1 I0 C0 F1}, diff_note={concepts: independent events | steps: probabilities of unions and intersections}, also_in={Star Review (Other) - Test 03 (04 Dec 2023) P2 Q1}]
$A$ and $B$ are two independent events of an experiment, such that $P(A)=\dfrac13$, $P(B)=\dfrac14$, then
\begin{qoptions}
\opt{A}{$P(A\cup B)=\dfrac12$}
\opt{B}{$P(A\cup\bar B)=\dfrac56$}
\opt{C}{$P(B\cup\bar A)=\dfrac34$}
\opt{D}{$P(A\cap B)=\dfrac1{12}$}
\end{qoptions}
\end{bankq}

\section*{TMAS papers}
\begin{bankq}[id=TMAS04-P2-MATH-S1-Q06, source={TMAS - Test 04 (19 Nov 2023)}, paper=2, qno=6, type=multi, subtopic={Bayes}, answer={B,D}, answer_src=key, difficulty=7, diff_parts={K1 S2 I2 C2 F1}, diff_note={concepts: conditional probability with priors | steps: prior proportional to k; probability of a biased coin; limit; Bayes}]
A box contains $n$ coins of which at least one is biased. Let $E_k$ denote the event that exactly $k$ out of the $n$ coins are biased. Let $P(E_k)$ is proportional to $k$ for $1\le k\le n$. If $A$ denotes the event that a coin selected out of $n$ coins at random, is biased, then :-
\begin{qoptions}
\opt{A}{$\displaystyle\lim_{n\to\infty}P(A)=\dfrac13$}
\opt{B}{$\displaystyle\lim_{n\to\infty}P(A)=\dfrac23$}
\opt{C}{$P\!\left(\dfrac{E_1}{A}\right)=\dfrac{3}{n(n+1)(2n+1)}$}
\opt{D}{$P\!\left(\dfrac{E_1}{A}\right)=\dfrac{6}{n(n+1)(2n+1)}$}
\end{qoptions}
\end{bankq}
\begin{bankq}[id=TMAS04-P2-MATH-S3-Q04, source={TMAS - Test 04 (19 Nov 2023)}, paper=2, qno=4, type=numerical, subtopic={Classical probability}, answer={3.00}, answer_src=key, difficulty=6, diff_parts={K1 S2 I2 C1 F1}, diff_note={concepts: probability from the minimum and maximum | steps: complement count; formula for the min; formula for the max; sum}]
Three numbers are chosen at random one after another with replacement, from the set $S=\{1,2,3,\ldots,100\}$. Let $P_1$ be the probability that the minimum of chosen numbers is at most $60$ and $P_2$ be the probability that the maximum of chosen numbers is at least $41$ and at most $80$. The value of $([P_1+P_2])$ is \underline{\hspace{0.8cm}}

(Where $[.]$ denotes greatest integer function)
\end{bankq}
\begin{bankq}[id=TMAS04-P2-MATH-S3-Q05, source={TMAS - Test 04 (19 Nov 2023)}, paper=2, qno=5, type=numerical, subtopic={Classical probability}, answer={1.775}, answer_src=key, difficulty=5, diff_parts={K1 S2 I1 C1 F1}, diff_note={concepts: inclusion exclusion probability | steps: multiples of 3 and 7; both; probability}]
A number is chosen at random from the set $\{1,2,3,\ldots,1000\}$. Let $P$ be the probability that the chosen number is a multiple of $3$ or a multiple of $7$. Then the value of $[5P]$ is \underline{\hspace{0.8cm}}

(Where $[.]$ denotes greatest integer function)
\end{bankq}

\section*{Star Review (Other) papers}
\begin{bankq}[id=STAR03-P1-MATH-S1-Q01, source={Star Review (Other) - Test 03 (04 Dec 2023)}, paper=1, qno=1, type=single, subtopic={Binomial distribution}, answer={B}, answer_src=key, difficulty=4, diff_parts={K1 S2 I1 C1 F0}, diff_note={concepts: probability of alternating colours with replacement | steps: two possible patterns; probability of each; add}]
A bag contains $5$ white and $3$ black balls. $4$ balls are drawn one by one, their color noted and replaced in the bag. Probability that white and black balls are drawn alternately is
\begin{qoptions}
\opt{A}{$\dfrac{15^2}{2^{10}}$}
\opt{B}{$\dfrac{15^2}{2^{11}}$}
\opt{C}{$\dfrac{15^2}{2^{12}}$}
\opt{D}{$\dfrac{15^2}{2^{13}}$}
\end{qoptions}
\end{bankq}
\begin{bankpassage}[id=STAR03-P1-MATH-S3-G1-2, source={Star Review (Other) - Test 03 (04 Dec 2023)}, questions={1-2}]
Let $A,B,C$ be $3$ exhaustive events of an experimental performance such that $P(A\cap B\cap C)=0$ and $P(A\cap B)=P(B\cap C)=P(C\cap A)=d\neq0$. If $P\left(A\bar B\cup\bar AB\right)=\dfrac{17}{24}$, $P\left(C\bar A\bar B\right)=\dfrac5{24}$ and $P\left(A\bar B\bar C\right)=\dfrac8{24}$. On the basis of above information, answer the following questions :
\end{bankpassage}
\begin{bankq}[id=STAR03-P1-MATH-S3-Q01, source={Star Review (Other) - Test 03 (04 Dec 2023)}, paper=1, qno=1, type=numerical, subtopic={Conditional probability}, answer={0.285}, answer_src=key, difficulty=7, diff_parts={K1 S2 I2 C2 F1}, diff_note={concepts: probabilities from union of events | steps: inclusion exclusion relations; find the common probabilities; conditional probability}, group=STAR03-P1-MATH-S3-G1-2]
The probability of occurence of event $A$ given that $B$ has occured and $C$ has not occured , is-
\end{bankq}
\begin{bankq}[id=STAR03-P1-MATH-S3-Q02, source={Star Review (Other) - Test 03 (04 Dec 2023)}, paper=1, qno=2, type=numerical, subtopic={Conditional probability}, answer={0.135}, answer_src=key, difficulty=6, diff_parts={K1 S2 I2 C1 F1}, diff_note={concepts: conditional probability with complementary event | steps: common probabilities; probability of the event not occurring; ratio}, group=STAR03-P1-MATH-S3-G1-2]
The probability of occurence of event $A$ \& $B$ given that $C$ has not occured, is-
\end{bankq}

\section*{Nurture (General) papers}
\begin{bankq}[id=NUR01-P1-MATH-S1-Q04, source={Nurture (General) - Test 01 (12 Jun 2021)}, paper=1, qno=4, type=multi, subtopic={Classical probability}, answer={D}, answer_src=solved, working={All row and column sums odd: 16 matrices out of 512 gives 1/32}, difficulty=7, diff_parts={K2 S2 I2 C1 F1}, diff_note={concepts: binary matrix with odd row and column sums, classical probability | steps: parity conditions; count matrices; probability}]
Let $A=[a_{ij}]_{3\times3}$ be a matrix where $a_{ij}\in\{0,1\}\ \forall\ 1\le i,j\le3$. Let $s_i=\sum\limits_{k=1}^{3}a_{ik}$ and $t_j=\sum\limits_{k=1}^{3}a_{kj}$ then the probability that the product $(s_1s_2s_3t_1t_2t_3)$ is an odd integer, is
\begin{qoptions}
\opt{A}{$\dfrac12$}
\opt{B}{$\dfrac18$}
\opt{C}{$\dfrac1{16}$}
\opt{D}{$\dfrac1{32}$}
\end{qoptions}
\end{bankq}
\begin{bankq}[id=NUR01-P2-MATH-S3-Q03, source={Nurture (General) - Test 01 (12 Jun 2021)}, paper=2, qno=3, type=single, subtopic={Classical probability}, answer={A}, answer_src=solved, working={20 points; pairs of equal parity classes 42 over C(20,2)=190: 21/95}, difficulty=6, diff_parts={K2 S2 I2 C1 F0}, diff_note={concepts: midpoint of two lattice points, classical probability | steps: parity classes; same class pairs; probability}]
The probability that mid point of the line segment of two points chosen belongs to $S$ is
\begin{qoptions}
\opt{A}{$\dfrac{21}{95}$}
\opt{B}{$\dfrac{42}{95}$}
\opt{C}{$\dfrac{36}{95}$}
\opt{D}{None of these}
\end{qoptions}
\end{bankq}
\begin{bankq}[id=NUR01-P2-MATH-S3-Q04, source={Nurture (General) - Test 01 (12 Jun 2021)}, paper=2, qno=4, type=single, subtopic={Classical probability}, answer={B}, answer_src=solved, working={16 favourable of 40 pairs on a vertical line: 2/5}, difficulty=6, diff_parts={K2 S2 I2 C1 F0}, diff_note={concepts: midpoint on a vertical line, classical probability | steps: count the pairs; parity of y; probability}]
If two points $(x_1,y_1)$ and $(x_2,y_2)$ are chosen such that $x_1=x_2$ and $y_1\neq y_2$ then probability that mid point of these two points belongs to $S$, is
\begin{qoptions}
\opt{A}{$\dfrac15$}
\opt{B}{$\dfrac25$}
\opt{C}{$\dfrac8{15}$}
\opt{D}{$\dfrac8{95}$}
\end{qoptions}
\end{bankq}
\end{document}
